\documentclass[10pt,a4paper]{amsart}
\usepackage[margin=19mm]{geometry}
\usepackage{amsmath,amssymb,mathtools}
\usepackage[scr=rsfs,cal=euler]{mathalfa}
\usepackage{xfrac}
\usepackage[unicode,hidelinks,bookmarks=false]{hyperref}
\newcommand{\ie}{i.e.\ }
\newcommand{\cf}{cf.\ }
\newcommand{\resp}{respectively\ }
\newcommand{\Subsection}{\subsection}
\providecommand{\nobreakdash}{\nobreak\hbox{-}}
\setlength{\emergencystretch}{2em}
\multlinegap0pt
\usepackage{amssymb}
\usepackage{xcolor}
\newtheorem{theorem}{Theorem}
\newtheorem{lemma}{Lemma}
\newtheorem{proposition}{Proposition}
\newcommand{\R}{\mathbb{R}}
\newcommand{\bal}{\begin{align}}
\newcommand{\bbal}{\begin{align*}}
\newcommand{\cd}{\cdot}
\newcommand{\dd}{\mathrm{d}}
\newcommand{\f}{\left}
\newcommand{\g}{\right}
\newcommand{\les}{\lesssim}
\newcommand{\na}{\nabla}
\newcommand{\pa}{\partial}
\title{On the continuous properties for the 3D incompressible rotating Euler equations}
\author{Jinlu Li, Yanghai Yu, Neng Zhu}
\date{Source: arXiv:2503.05069v2 (CC BY 4.0)}
\begin{document}
\maketitle

\noindent\emph{Source and licence.} On the continuous properties for the 3D incompressible rotating Euler equations. Version arXiv:2503.05069v2 (CC BY 4.0), distributed by arXiv under the Creative Commons Attribution 4.0 International licence, http://creativecommons.org/licenses/by/4.0/, as recorded in the arXiv licence metadata for this exact version and shown on the abstract page https://arxiv.org/abs/2503.05069v2. Full licence notice: \texttt{licenses/arxiv-2503.05069-CC-BY-4.0.txt}.\par\smallskip

\noindent\textit{Editorial scope.} Counted unit: the article's proposition pro1 (source0.tex lines 486--497). The article's complete original proof of it is delivered with it (source0.tex lines 498--538, one \texttt{proof} environment of 41 lines). No mathematics is written, repaired or re-derived: the statement and the proof are the article's own lines, in the article's own notation, and no retained line is substituted or re-flowed.  Retained uncounted as the source-local context the proof needs: lines 94--100; lines 175--185; lines 177--179; lines 199--209; lines 357--367.  Nothing was omitted from any retained span, so the package carries no omission record.  The retained lines cite 1 works, whose entries are transcribed verbatim from the article's own bibliography source in the e-print and delivered with short editorial labels recorded in the item's reference mapping.  No retained line contains a figure environment, an image-inclusion command or drawing code, so no image is delivered and the package declares no asset dependency.  Editorial formatting only, in addition: an AMS \texttt{amsart} preamble that restates the article's own declarations and macros used by the retained lines, namely the environments \texttt{theorem}, \texttt{lemma}, \texttt{proposition} on separate counters, together with the standard packages those lines need.  The retained environments print in this document as Theorem 1, Lemma 1, Proposition 1 in order of appearance, whereas the article prints the same items with the numbering of its own full document.  The complete native source of the article as distributed in the arXiv e-print for this exact version is bundled unchanged as \texttt{sources/174360/source0.tex}.
\begin{align}\label{CE}
\begin{cases}
\pa_t u+\Omega e_3\times u+u\cdot \nabla u+\nabla P=0, &\quad (t,x)\in \R^+\times\R^3,\\
\mathrm{div\,} u=0,&\quad (t,x)\in \R^+\times\R^3,\\
u(t=0,x)=u_0(x), &\quad x\in \R^3,
\end{cases}
\end{align}

\begin{theorem}\label{th1}
Assume that $u_0\in B^s_{p,r}(\R^3)$ with $(s,p,r)$ satisfying
\begin{align}\label{eq:spr1}
s>\frac{3}{p}+1, (p,r)\in (1,\infty)\times [1,\infty] \quad   \mathrm{or}    \quad s=\frac{3}{p}+1, (p,r)\in (1,\infty)\times\{1\}.
\end{align}
Then there exists a short time $T>0$ such that the Euler-Coriolis system \eqref{CE} admits a unique solution $u\in L^\infty([0,T];B^s_{p,r})$ satisfying that
\begin{align}\label{hx}
\|u(t)\|_{L^\infty_TB^s_{p,r}}\les\|u_0\|_{B^s_{p,r}}.
\end{align}
Furthermore, if $r<\infty$, the above solution $u\in C([0,T];B^s_{p,r})$ is continuous dependent in $B^s_{p,r}$.
\end{theorem}

\begin{align}
s>\frac{3}{p}+1, (p,r)\in (1,\infty)\times [1,\infty] \quad   \mathrm{or}    \quad s=\frac{3}{p}+1, (p,r)\in (1,\infty)\times\{1\}.
\end{align}

\begin{theorem}\label{th2}
Assume that $(s,p,r)$ satisfying \eqref{eq:spr1}.  The data-to-solution map $u_0\mapsto \mathbf{S}_{t}(u_0)$ of the Cauchy problem \eqref{CE} is not uniformly continuous from any bounded subset in $U_R\subset B^s_{p,r}$ into $\mathcal{C}([0,T];B^s_{p,r})$. More precisely, there exist two initial data sequences $f_n+g_n$ and $f_n$ such that
\bbal
&\f\|(f_n)^{(2)}\g\|_{B^s_{p,r}}\approx 1 \quad\text{and}\\
 &\lim_{n\rightarrow \infty}\f\|(f_n)^{(1)}\g\|_{B^s_{p,r}}=\lim_{n\rightarrow \infty}\f\|(f_n)^{(3)}\g\|_{B^s_{p,r}}=\lim_{n\rightarrow \infty}\|g_n\|_{B^s_{p,r}}= 0,
\end{align*}
but the distance of second component of the difference between the solution sequences $\mathbf{S}_t(f_n+g_n)$ and $\mathbf{S}_t(f_n)$
\bbal
\liminf_{n\rightarrow \infty}\f\|\f(\mathbf{S}_t(f_n+g_n)-\mathbf{S}_t(f_n)\g)^{(2)}\g\|_{B^s_{p,r}}\gtrsim t,  \quad \forall \;t\in[0,T].
\end{align*}
\end{theorem}

\begin{lemma}[\cite{B}]\label{lp}
Assume $(s,p,r)$ satisfies \eqref{eq:spr1} and $\sigma>0$. Then
 there exists a constant $C$, depending only on $d,p,r,\sigma$ or $s$ such that
$$\|fg\|_{B^{\sigma}_{p,r}}\leq C\f(\|f\|_{L^\infty}\|g\|_{B^\sigma_{p,r}}+\|g\|_{L^\infty}\|f\|_{B^{\sigma}_{p,r}}\g),\quad \forall f,g\in L^\infty\cap B^{\sigma}_{p,r}.$$
Furthermore, due to the embedding $B^{s-1}_{p,r}(\R^3)\hookrightarrow L^{\infty}(\R^3)$, there holds
$$\|f\cdot \na g\|_{B^{s-1}_{p,r}}\leq C\|f\|_{B^{s-1}_{p,r}}\|g\|_{B^s_{p,r}},\quad \forall(f,g)\in B^{s-1}_{p,r}\times B^s_{p,r}$$
and
\begin{align}\label{cj}
\|f\cdot \na g\|_{B^{s}_{p,r}}\leq C\f(\|f\|_{B^{s-1}_{p,r}}\|g\|_{B^{s+1}_{p,r}}+\|f\|_{B^{s}_{p,r}}\|g\|_{B^{s}_{p,r}}\g),\quad \forall(f,g)\in B^{s}_{p,r}\times B^{s+1}_{p,r}.
\end{align}
\end{lemma}

\begin{proposition}\label{pro1}
Under the assumptions of Theorem \ref{th2}, we have
\bal\label{et0}
\|\mathbf{S}_{t}(u_0)-u_0+t\mathbf{v}_0(u_0)\|_{B^{s}_{p,r}}\leq Ct^{2}\mathbf{E}(u_0),
\end{align}
here and in what follows we denote
\bbal
\mathbf{E}(u_0)&:=1+\|u_0\|_{B^{s-1}_{p,r}}\f(\|u_0\|_{B^{s+1}_{p,r}}+
\|u_0\|_{B^{s-1}_{p,r}}
\|u_0\|_{B^{s+2}_{p,r}}\g).
\end{align*}
\end{proposition}

\begin{proof}\,
For simplicity, we denote $u(t)=\mathbf{S}_t(u_0)$ here and in what follows. By Theorem \ref{th1}, one has
\bal\label{bound:u}
\|u(t)\|_{B^{s+k}_{p,r}}\leq C\|u_0\|_{B^{s+k}_{p,r}}\quad \text{for }\; k\in\{0,\pm1,2\}.
\end{align}
Using the Mean Value Theorem and Lemma \ref{lp}, we obtain from \eqref{CE} that
\bal\label{et1}
&\|u(t)-u_0\|_{B^{s-1}_{p,r}}\leq\int^t_0\|\pa_\tau u\|_{B^{s-1}_{p,r}} \dd\tau
\nonumber\\
\leq&~ \int^t_0\f(\|u\cd\na u\|_{B^{s-1}_{p,r}}+\|\mathcal{P}(\Omega e_3\times u)\|_{B^{s-1}_{p,r}}+\|\mathcal{Q}(u\cd \na u)\|_{B^{s-1}_{p,r}}\g)\dd \tau
\nonumber\\
\leq&~  C\int^t_0\f(\|u\|_{B^s_{p,r}}\|u\|_{B^{s-1}_{p,r}}+\|u\|_{B^{s-1}_{p,r}}\g)\dd \tau\nonumber\\
\leq&~  Ct\|u_0\|_{B^{s-1}_{p,r}},
\end{align}
\bal\label{et2}
&\|u(t)-u_0\|_{B^{s}_{p,r}}\leq\int^t_0\|\pa_\tau u\|_{B^{s}_{p,r}} \dd\tau\nonumber\\
\leq&~ \int^t_0\f(\|u\cd\na u\|_{B^{s}_{p,r}}+\|\mathcal{P}(\Omega e_3\times u)\|_{B^{s}_{p,r}}+\|\mathcal{Q}(u\cd \na u)\|_{B^{s}_{p,r}}\g)\dd \tau\nonumber\\
\leq&~ Ct\f(1+\|u_0\|_{B^{s-1}_{p,r}}\|u_0\|_{B^{s+1}_{p,r}}\g),
\end{align}
\bal\label{et3}
&\|u(t)-u_0\|_{B^{s+1}_{p,r}}\leq\int^t_0\|\pa_\tau u\|_{B^{s+1}_{p,r}} \dd\tau
\nonumber\\
\leq&~ \int^t_0\f(\|u\cd\na u\|_{B^{s+1}_{p,r}}+\|\mathcal{P}(\Omega e_3\times u)\|_{B^{s+1}_{p,r}}+\|\mathcal{Q}(u\cd \na u)\|_{B^{s+1}_{p,r}}\g)\dd \tau
\nonumber\\
\leq&~ Ct\f(\|u_0\|_{B^{s+1}_{p,r}}+\|u_0\|_{B^{s-1}_{p,r}}\|u_0\|_{B^{s+2}_{p,r}}\g).
\end{align}
Using the Mean Value Theorem once again, we obtain that
\bal\label{et4}
&\|u(t)-u_0+t\mathbf{v}_0(u_0)\|_{B^{s}_{p,r}}\leq \int^t_0\|\pa_\tau u+\mathbf{v}_0(u_0)\|_{B^{s}_{p,r}} \dd\tau
\nonumber\\
\leq&~\int^t_0\|\mathcal{Q}(u)-\mathcal{Q}(u_0)\|_{B^{s}_{p,r}} \dd\tau+\int^t_0\|u\cd\na u-u_0\cd\na u_0\|_{B^{s}_{p,r}} \dd\tau\nonumber\\
\quad&+\Omega\int^t_0\|\mathcal{P}\f(e_3\times (u-u_0)\g)\|_{B^{s}_{p,r}}\dd \tau
\nonumber\\
\lesssim&~ \int^t_0\|u(\tau)-u_0\|_{B^{s}_{p,r}} \dd\tau+\int^t_0\|u(\tau)-u_0\|_{L^\infty} \|u(\tau)\|_{B^{s+1}_{p,r}} \dd\tau
\nonumber\\
\quad & + \int^t_0\|u(\tau)-u_0\|_{B^{s+1}_{p,r}}  \|u_0\|_{L^\infty}\dd \tau\nonumber\\
\lesssim&~ \int^t_0\|u(\tau)-u_0\|_{B^{s}_{p,r}} \dd\tau+\|u_0\|_{B^{s+1}_{p,r}}\int^t_0\|u(\tau)-u_0\|_{L^\infty}  \dd\tau\nonumber\\
\quad &+ \|u_0\|_{L^\infty}\int^t_0\|u(\tau)-u_0\|_{B^{s+1}_{p,r}}  \dd \tau.
\end{align}
Plugging \eqref{et1}-\eqref{et3} into \eqref{et4} and using the fact that $B^{s-1}_{p,r}(\R^3)\hookrightarrow L^\infty(\R^3)$, yields the desired result of Proposition \ref{pro1}.
\end{proof}

\begin{thebibliography}{99}
\bibitem[B]{B} H. Bahouri, J. Y. Chemin, R. Danchin, Fourier Analysis and Nonlinear Partial Differential Equations, Grundlehren der Mathematischen Wissenschaften, Springer, Heidelberg, 2011.
\end{thebibliography}
\end{document}
